\section*{Limitations}

\paragraph{Annotator disagreement and label reliability.}
The labels themselves are not fully reliable.
The dataset's annotators agree only moderately, with a Fleiss' $\kappa$ of 0.644 for the clarity level and 0.48 for the nine evasion types \citep{thomas-etal-2024-never}, and the hardest pairs reach only 0.43 (\emph{General} vs.\ \emph{Implicit}) and 0.56 (\emph{General} vs.\ \emph{Deflection}) \citep{thomas-etal-2026-semeval}.
On our development set only 125 of the 308 items are unanimous, while 150 split two to one and 33 receive three different labels (\S\ref{sec:task}).
This affects both training and evaluation.
Every training item carries a single label from one annotator, so the models learn one reading of answers that other annotators would label differently, and they cannot learn how disputed an item is.
Multi-reference scoring credits a prediction that matches \emph{any} annotator, so a ``correct'' prediction may be one that most annotators rejected.
Scores also depend on how many annotators an item has, and rescored with two annotators, as on the test set, our 8B system drops from 0.543 to 0.524.
The task has a low ceiling as well, since a human annotator scored against the other two reaches only 0.643--0.766 (mean 0.684), so differences near that level fall within human disagreement.
A predicted label such as \emph{Dodging} is therefore one plausible reading of an answer, not a fact about it or a verdict on the speaker, and it should be reported with its uncertainty.

\paragraph{Development-set evaluation.}
All results are on the 308-item development set, because the test labels were never released and the evaluation server has closed.
The development set has three annotators per item and the test set two.
We rescore with every two-annotator subset to approximate the test setting, but we cannot measure the difference in answer length between the splits.
With only 308 items, differences between systems carry wide intervals.
The gain of the 8B system has a 95\% bootstrap interval of [0.054, 0.224], and the all-data variant could not be separated from its control.

\paragraph{Seeds and selection.}
Seed variance is large (the standard deviation of single-model Subtask~2 scores is 0.043 for the 8B classifier), which is why our claims rest on ten paired seeds.
The twelve-epoch schedule, however, was screened on three seeds only, so it is adopted but not yet established.
Epochs are chosen on a 345-row slice of the training data, which is itself noisy.
The one decision-rule parameter is fitted by nested cross-validation on the development set, which keeps the score held out but means that the system uses development labels to fit one scalar.

\paragraph{Scope of the comparison.}
Changing the backbone also changes the training recipe that comes with it (learning rate, epochs and parameter-efficient fine-tuning), so ``knowledge'' here is shorthand for everything a larger pretrained model brings.
We cannot rule out that the pretraining data of the 8B model includes the public interview transcripts the dataset is drawn from.
The published comparison points come from other teams' papers and protocols, which we did not control.

\paragraph{Coverage.}
None of our models ever predicts the rarest class (\emph{Partial/half-answer}, 14 development reference sets), and the study covers one dataset in one language (English, U.S.\ presidential interviews).
